\documentclass[12pt,a4paper]{article}

% Packages
\usepackage[utf8]{inputenc}
\usepackage[T1]{fontenc}
\usepackage{graphicx}
\usepackage{amsmath}
\usepackage{amssymb}
\usepackage{hyperref}
\usepackage[margin=1in]{geometry}
\usepackage[backend=biber,style=numeric,sorting=none]{biblatex}
\addbibresource{biblio.bib}


% Title and authors
\title{Multi-View Neural Network for Musculoskeletal Disorder Detection in Radiographic Imaging}
\author{
    Camille Trouvé\textsuperscript{1},
    Ceren Dinc\textsuperscript{1},
    Esteban Boulanger\textsuperscript{1} \\
    \textsuperscript{1}\textit{Faculté de Gestion Economie et Sciences - Ecole du Numérique} \\
}
\date{\today}

\begin{document}

\maketitle

% Abstract
\begin{abstract}
Musculoskeletal disorders are among the main causes of disability, according to the World Health Organization. This study aims to develop a tool capable of determining abnormalities in musculoskeletal radiographs. The proposed model is a transformer-based deep neural network with a multiple-view fusion module to integrate multiple radiographic projections, mimicking how radiologists interpret radiographic views. The model was trained and tested on the MURA dataset (Stanford), a large open-access dataset combining 40,561 radiographic images of the upper extremity from multiple studies. The model achieved an accuracy of XX (to be determined). The results are expected to demonstrate the benefit of multi-view integration for improving diagnostic performance and robustness in musculoskeletal abnormality detection. 
\end{abstract}

% Introduction
\section{Introduction}
\label{sec:introduction}

Write your introduction here.

% Material and Method
\section{Material and Method}
\label{sec:material-method}

Write your material and method section here.

% Results
\section{Results}
\label{sec:results}

Write your results here.

% Discussion
\section{Discussion}
\label{sec:discussion}

Write your discussion here.

\printbibliography

\end{document}
